% GENERATED FILE. Edit canonical lessons, then run npm run generate:books.
% canonical-source-hash: fnv1a64:915d3bb49074a390
% canonical-lessons: MR-C462-heva, MR-C462-dves, MR-C462-kilas, MR-C462-moh, MR-C462-lobh

\chapter{Envy and Greed}
\label{ch:mr-envy-and-greed}
\hlchaptermodality{462}

\begin{chapteropening}
\textbf{By the end of this chapter:} \emph{I can talk about envy, hatred, disgust, temptation and greed.}

Complete the last lesson of chapter 462: I can talk about envy, hatred, disgust, temptation and greed.
\end{chapteropening}

\section[\texorpdfstring{\mr{हेवा} (hevā)}{हेवा (hevā)}]{\mr{हेवा} (hevā) — envy}

\label{lesson:MR-C462-heva}



\begin{quote}
Before the new one: say the Marathi for to hurt someone's feelings, then the Marathi for to belittle.
\end{quote}

\subsection*{You'll want to know: \mr{हेवा}}
\textbf{\mr{हेवा}} — \emph{hevā} — \textquotedblleft{}envy\textquotedblright{}.

\begin{quote}
\mr{मित्राच्या} \mr{नव्या} \mr{सायकलचा} \mr{त्याला} \mr{हेवा} \mr{वाटला}.
\end{quote}

\begin{quote}
\emph{mitrācyā navyā sāykalcā tyālā hevā vāṭlā.}
\end{quote}

\begin{quote}
He envied his friend's new bicycle.
\end{quote}

Masculine.

\mr{हेवा} \mr{वाटणे}, to envy: the dative for the one who envies, the genitive for the thing. Milder than \mr{मत्सर}.

\begin{cousinweb}
\textbf{\mr{हेवा}} (\emph{hevā}) is a Marathi noun whose origin is not settled.
\end{cousinweb}

\begin{grammarlens}[title={What you've built}]
One more word for describing what you lived and felt.
\end{grammarlens}

\subsection*{Your turn}
\begin{itemize}
\raggedright
  \item \emph{Say it:} \emph{hevā}
  \item \emph{Say it:} \emph{hevā}, once more
  \item \emph{Say it:} \emph{mitrācyā navyā sāykalcā tyālā hevā vāṭlā.} — He envied his friend's new bicycle.
  \item \emph{Recall:} say the Marathi for to hurt someone's feelings, then the Marathi for to belittle, then say \emph{hevā} again
\end{itemize}

\begin{tcolorbox}[breakable,colback=teal!4,colframe=teal!35!black,title={Before you move on}]
What does \mr{हेवा} mean? (\textquotedblleft{}Envy\textquotedblright{}.) Say it once more.
\end{tcolorbox}
\section[\texorpdfstring{\mr{द्वेष} (dveṣ)}{द्वेष (dveṣ)}]{\mr{द्वेष} (dveṣ) — hatred}

\label{lesson:MR-C462-dves}



\begin{quote}
Before the new one: say the Marathi for to belittle, then the Marathi for envy.
\end{quote}

\subsection*{You'll want to know: \mr{द्वेष}}
\textbf{\mr{द्वेष}} — \emph{dveṣ} — \textquotedblleft{}hatred\textquotedblright{}.

\begin{quote}
\mr{मनात} \mr{द्वेष} \mr{ठेवू} \mr{नये}.
\end{quote}

\begin{quote}
\emph{manāt dveṣ ṭhevū naye.}
\end{quote}

\begin{quote}
One should not keep hatred in one's heart.
\end{quote}

Masculine.

\mr{द्वेष} \mr{करणे}, to hate, with the genitive: \mr{त्याचा} \mr{द्वेष} \mr{करणे}.

\begin{cousinweb}
\textbf{\mr{द्वेष}} (\emph{dveṣ}) is a tatsama, taken from Sanskrit \emph{dveṣa}, \textquotedblleft{}hatred\textquotedblright{} (CDIAL 6696).
\end{cousinweb}

\begin{grammarlens}[title={What you've built}]
One more word for describing what you lived and felt.
\end{grammarlens}

\subsection*{Your turn}
\begin{itemize}
\raggedright
  \item \emph{Say it:} \emph{dveṣ}
  \item \emph{Say it:} \emph{dveṣ}, once more
  \item \emph{Say it:} \emph{manāt dveṣ ṭhevū naye.} — One should not keep hatred in one's heart.
  \item \emph{Recall:} say the Marathi for to belittle, then the Marathi for envy, then say \emph{dveṣ} again
\end{itemize}

\begin{tcolorbox}[breakable,colback=teal!4,colframe=teal!35!black,title={Before you move on}]
What does \mr{द्वेष} mean? (\textquotedblleft{}Hatred\textquotedblright{}.) Say it once more.
\end{tcolorbox}
\section[\texorpdfstring{\mr{किळस} (kiḷas)}{किळस (kiḷas)}]{\mr{किळस} (kiḷas) — disgust}

\label{lesson:MR-C462-kilas}



\begin{quote}
Before the new one: say the Marathi for envy, then the Marathi for hatred.
\end{quote}

\subsection*{You'll want to know: \mr{किळस}}
\textbf{\mr{किळस}} — \emph{kiḷas} — \textquotedblleft{}disgust\textquotedblright{}.

\begin{quote}
\mr{घाणेरडं} \mr{पाणी} \mr{पाहून} \mr{मला} \mr{किळस} \mr{आली}.
\end{quote}

\begin{quote}
\emph{ghāṇerḍaṁ pāṇī pāhūn malā kiḷas ālī.}
\end{quote}

\begin{quote}
Seeing the filthy water, I felt disgust.
\end{quote}

Feminine.

\mr{किळस} \mr{येणे}, to feel disgust, with the dative; \mr{किळसवाणा}, disgusting.

\begin{cousinweb}
\textbf{\mr{किळस}} (\emph{kiḷas}) is inherited: CDIAL lists the Marathi verb \emph{kiḷasṇẽ}, \textquotedblleft{}to loathe\textquotedblright{}, under Sanskrit \emph{kliśyate}, \textquotedblleft{}is troubled\textquotedblright{} (CDIAL 3623).
\end{cousinweb}

\begin{grammarlens}[title={What you've built}]
One more word for describing what you lived and felt.
\end{grammarlens}

\subsection*{Your turn}
\begin{itemize}
\raggedright
  \item \emph{Say it:} \emph{kiḷas}
  \item \emph{Say it:} \emph{kiḷas}, once more
  \item \emph{Say it:} \emph{ghāṇerḍaṁ pāṇī pāhūn malā kiḷas ālī.} — Seeing the filthy water, I felt disgust.
  \item \emph{Recall:} say the Marathi for envy, then the Marathi for hatred, then say \emph{kiḷas} again
\end{itemize}

\begin{tcolorbox}[breakable,colback=teal!4,colframe=teal!35!black,title={Before you move on}]
What does \mr{किळस} mean? (\textquotedblleft{}Disgust\textquotedblright{}.) Say it once more.
\end{tcolorbox}
\section[\texorpdfstring{\mr{मोह} (moh)}{मोह (moh)}]{\mr{मोह} (moh) — temptation}

\label{lesson:MR-C462-moh}



\begin{quote}
Before the new one: say the Marathi for hatred, then the Marathi for disgust.
\end{quote}

\subsection*{You'll want to know: \mr{मोह}}
\textbf{\mr{मोह}} — \emph{moh} — \textquotedblleft{}temptation\textquotedblright{}.

\begin{quote}
\mr{लाडू} \mr{पाहून} \mr{मला} \mr{खाण्याचा} \mr{मोह} \mr{झाला}.
\end{quote}

\begin{quote}
\emph{lāḍū pāhūn malā khāṇyācā moh jhālā.}
\end{quote}

\begin{quote}
Seeing the ladoos, I was tempted to eat one.
\end{quote}

Masculine.

\mr{मोह} \mr{होणे}, to be tempted, with the dative; \mr{मोह} \mr{टाळणे}, to resist temptation.

\begin{cousinweb}
\textbf{\mr{मोह}} (\emph{moh}) is a tatsama, taken from Sanskrit \emph{moha}, \textquotedblleft{}bewilderment\textquotedblright{} (CDIAL 10360); in Marathi it is the pull of temptation.
\end{cousinweb}

\begin{grammarlens}[title={What you've built}]
One more word for describing what you lived and felt.
\end{grammarlens}

\subsection*{Your turn}
\begin{itemize}
\raggedright
  \item \emph{Say it:} \emph{moh}
  \item \emph{Say it:} \emph{moh}, once more
  \item \emph{Say it:} \emph{lāḍū pāhūn malā khāṇyācā moh jhālā.} — Seeing the ladoos, I was tempted to eat one.
  \item \emph{Recall:} say the Marathi for hatred, then the Marathi for disgust, then say \emph{moh} again
\end{itemize}

\begin{tcolorbox}[breakable,colback=teal!4,colframe=teal!35!black,title={Before you move on}]
What does \mr{मोह} mean? (\textquotedblleft{}Temptation\textquotedblright{}.) Say it once more.
\end{tcolorbox}
\section[\texorpdfstring{\mr{लोभ} (lobh)}{लोभ (lobh)}]{\mr{लोभ} (lobh) — greed}

\label{lesson:MR-C462-lobh}



\begin{quote}
Before the new one: say the Marathi for disgust, then the Marathi for temptation.
\end{quote}

\subsection*{You'll want to know: \mr{लोभ}}
\textbf{\mr{लोभ}} — \emph{lobh} — \textquotedblleft{}greed\textquotedblright{}.

\begin{quote}
\mr{पैशाच्या} \mr{लोभाने} \mr{त्याने} \mr{मित्राला} \mr{फसवलं}.
\end{quote}

\begin{quote}
\emph{paiśācyā lobhāne tyāne mitrālā phasavlaṁ.}
\end{quote}

\begin{quote}
Out of greed for money he cheated his friend.
\end{quote}

Masculine.

\mr{लोभी}, greedy.

\begin{cousinweb}
\textbf{\mr{लोभ}} (\emph{lobh}) is a tatsama, taken from Sanskrit \emph{lobha}, \textquotedblleft{}eager desire\textquotedblright{} (CDIAL 11147).
\end{cousinweb}

\begin{grammarlens}[title={What you've built}]
One more word for describing what you lived and felt.
\end{grammarlens}

\subsection*{Your turn}
\begin{itemize}
\raggedright
  \item \emph{Say it:} \emph{lobh}
  \item \emph{Say it:} \emph{lobh}, once more
  \item \emph{Say it:} \emph{paiśācyā lobhāne tyāne mitrālā phasavlaṁ.} — Out of greed for money he cheated his friend.
  \item \emph{Recall:} say the Marathi for disgust, then the Marathi for temptation, then say \emph{lobh} again
\end{itemize}

\begin{tcolorbox}[breakable,colback=teal!4,colframe=teal!35!black,title={Before you move on}]
What does \mr{लोभ} mean? (\textquotedblleft{}Greed\textquotedblright{}.) Say it once more.
\end{tcolorbox}
